\section{Nonintegrable phase factors (11/18)}

\subsection{A review of adiabatic evolution and Berry's phase}
Let
\begin{equation}
 \hat H(\lambda)\dket{n,\lambda}
 =E_n(\lambda)\dket{n,\lambda}.
\end{equation}
In the adiabatic limit,
\begin{equation}
 \dket{\psi(t)}
 =\sum_n C_n(t)
 \exp\left[-\frac{\imth}{\hbar}
 \int_0^t E_n(t')\dif t'\right]
 \dket{n,t},
\end{equation}
where
\begin{equation}
 \dot C_n(t)
 =-C_n(t)\dbra{n,t}\frac{\dif}{\dif t}\dket{n,t}.
\end{equation}
Thus
\begin{equation}
 C_n(t)=C_n(0)
 \exp\left[-\imth\int_{\lambda(0)}^{\lambda(t)}
 \mathcal A_n(\lambda)\dif\lambda\right],
\end{equation}
with the \emph{Berry connection} defined as
\begin{equation}
 \boxed{
 \mathcal A_n(\lambda)
 =-\imth\dbra{n,\lambda}\partial_\lambda\dket{n,\lambda}.}
\end{equation}
Since $\expval{n,\lambda|n,\lambda}=1$,
\begin{equation}
 \mathcal A_n^*(\lambda)=\mathcal A_n(\lambda).
\end{equation}
For several slowly-varying parameters $\{\lambda_j\}$, there can be more than one component of Berry connection 
\begin{equation}
 \mathcal A_{n,j}
 =-\imth\dbra n\partial_{\lambda_j}\dket n.
\end{equation}
and 
\begin{equation}
 C_n(t)=C_n(0)
 \exp\left[-\imth\sum_j\int_{\lambda(0)}^{\lambda(t)}
 \mathcal A_{n,j}(\lambda)\dif\lambda_j\right],
\end{equation}
Notice that the phase factor in $C_n(t)$ depends not only on the endpoints $\lambda(0),\lambda(t)$, but also on the full trajectory in parameter space. Therefore it is a non-integrable phase factor. 

\subsection{Gauge redundancy}
Under
\begin{equation}
 \dket{n,\lambda}
 \longmapsto
 \dket{\widetilde n,\lambda}
 =e^{\imth f(\lambda)}\dket{n,\lambda},
\end{equation}
the Berry connection transforms as
\begin{equation}
 \boxed{
 \widetilde{\mathcal A}_n(\lambda)
 =\mathcal A_n(\lambda)+\frac{\dif f}{\dif\lambda}.}
\end{equation}
For an open path,
\begin{equation}
 e^{\imth\gamma_n}
 =\exp\left[-\imth\int_{\lambda_0}^{\lambda_1}
 \mathcal A_n(\lambda)\dif\lambda\right]
\end{equation}
changes by an endpoint phase difference $f(\lambda_1)-f(\lambda_0)$. For a closed contour $C$,
\begin{equation}
 \boxed{
 e^{\imth\gamma_n(C)}
 =\exp\left[-\imth\oint_C
 \mathcal A_n(\lambda)\dif\lambda\right]}
\end{equation}
is gauge invariant.

\subsection{Aharonov--Bohm phase}

Consider
\begin{equation}
 \hat H(\bm\lambda)
 =\frac{(\hat{\bm p}-e\bm A)^2}{2m}
 +e\phi(\hat{\bm r})
 +V(\hat{\bm r}-\bm\lambda),
\end{equation}
where $V$ traps the particle near $\bm\lambda$. For a slow displacement
$\bm\lambda\to\bm\lambda+\bm\delta$,
\begin{equation}
 \dket{n,\bm\lambda+\bm\delta}
 \simeq
 e^{-\imth\hat{\bm\pi}\cdot\bm\delta/\hbar}
 \dket{n,\bm\lambda},
 \qquad
 \hat{\bm\pi}=\hat{\bm p}-e\bm A.
\end{equation}
Hence
\begin{equation}
 -\imth\dbra{n,\bm\lambda}
 \frac{\partial}{\partial\lambda_j}
 \dket{n,\bm\lambda}
 =\frac1\hbar\dbra n\hat\pi_j\dket n
 \simeq-\frac{e}{\hbar}A_j(\bm\lambda).
\end{equation}
The geometric phase is therefore
\begin{equation}
 \boxed{
 e^{\imth\gamma}
 =\exp\left[
 \frac{\imth e}{\hbar}
 \int_{\bm\lambda(0)}^{\bm\lambda(t)}
 \bm A(\bm r)\cdot\dif\bm r\right].}
\end{equation}
Under
\begin{equation}
 \bm A\longmapsto\bm A-\bm\nabla f,
 \qquad
 \psi\longmapsto e^{-\imth ef/\hbar}\psi,
\end{equation}
the induced endpoint phase agrees with the gauge transformation of the wave
function.


For a closed contour $C$,
\begin{equation}
 e^{\imth\gamma_{\mathrm{AB}}}
 =\exp\left[
 \frac{\imth e}{\hbar}\oint_C\bm A\cdot\dif\bm r\right]
 =\exp\left(2\pi\imth\frac{\Phi}{\Phi_0}\right),
\end{equation}
where
\begin{equation}
 \boxed{\Phi_0=\frac{2\pi\hbar}{e}=\frac{h}{e}}
\end{equation}
is the flux quantum.

\subsection{Particle on a ring}
For a particle on a ring of radius $R$ threaded by flux $\Phi$,
\begin{equation}
 \hat H
 =\frac{(\hat p-eA)^2}{2m}
 =\frac{\hbar^2}{2mR^2}
 \left(-\imth\partial_\theta-\phi\right)^2,
\end{equation}
where
\begin{equation}
 \phi=\frac{eAR}{\hbar}
 =\frac{\Phi}{\Phi_0}.
\end{equation}
Periodic boundary conditions give
\begin{equation}
 \psi_n(\theta)=\frac{e^{\imth n\theta}}{\sqrt{2\pi}},
 \qquad n\in\mathbb Z,
\end{equation}
and
\begin{equation}
 \boxed{
 E_n(\phi)=\frac{\hbar^2}{2mR^2}(n-\phi)^2.}
\end{equation}
Therefore,
\begin{equation}
 E_{n+1}(\phi+1)=E_n(\phi),
\end{equation}
so the spectrum is periodic in $\Phi$ with period $\Phi_0$.

For two spinless electrons and $0\leq\phi\leq1$, the total energy is given by
\begin{equation}
 U(\phi)=E_0(\phi)+E_1(\phi)
 =\frac{\hbar^2}{2mR^2}
 \left[\phi^2+(\phi-1)^2\right].
\end{equation}
The magnetization is
\begin{equation}
 \boxed{
 M=-\frac{\partial U}{\partial\Phi}
 =\frac{\hbar^2}{mR^2\Phi_0}(1-2\phi).}
\end{equation}
It is a periodic function of $\phi$ with a periodicity of 1. 

\subsection{Dirac quantization}
Cover a sphere surrounding a magnetic monopole by northern and southern
patches $R$ and $\overline R$, with common boundary $C$. Stokes' theorem gives
\begin{equation}
 \oint_C\bm A_R\cdot\dif\bm r
 =\int_R\bm B\cdot\dif\bm a,
\end{equation}
and
\begin{equation}
 \oint_C\bm A_{\overline R}\cdot\dif\bm r
 =-\int_{\overline R}\bm B\cdot\dif\bm a.
\end{equation}
Single-valuedness of the phase requires
\begin{equation}
 \frac{e}{\hbar}
 \left(\Phi_R+\Phi_{\overline R}\right)
 =2\pi N,
 \qquad N\in\mathbb Z.
\end{equation}
Thus the total monopole flux is quantized:
\begin{equation}
 \boxed{
 \Phi_{S^2}=N\Phi_0,
 \qquad N\in\mathbb Z.}
\end{equation}